\documentclass[11pt]{article}
\usepackage[margin=1in]{geometry}
\usepackage{amsmath, amssymb, amsfonts}
\usepackage{booktabs}
\usepackage{siunitx}
\usepackage{enumitem}
\usepackage{physics}
\usepackage{microtype}
\usepackage{hyperref}
\hypersetup{colorlinks=true, linkcolor=blue, urlcolor=blue, citecolor=blue}

\setlist[enumerate]{leftmargin=*, itemsep=0.25em}
\sisetup{round-mode=places,round-precision=2}

\newcommand{\eq}{\Leftrightarrow}

\title{Tutorial 1}
\author{ECON 101 -- Macroeconomics}
\date{Fall 2025}

\begin{document}
\maketitle



%==========================
% Q1
%==========================
\noindent \textbf{Q1. Market Equilibrium and Price Controls}

Demand and supply for coffee are:
\[
Q_d=100-2P,\qquad Q_s=20+3P.
\]
\begin{enumerate}[label=(\alph*)]
\item Find the equilibrium price and quantity.
\item With a price floor at \(P=20\), compute the surplus.
\item With a price ceiling at \(P=10\), compute the shortage.
\end{enumerate}

\textbf{Solution.}
\begin{enumerate}[label=(\alph*)]
\item Equilibrium: set \(Q_d=Q_s\).
\[
100-2P=20+3P \;\eq\; 80=5P \;\eq\; P^\ast=16,\quad Q^\ast=100-2(16)=68.
\]
\item At \(P=20\): \(Q_d=100-40=60,\; Q_s=20+60=80\).
Surplus \(=Q_s-Q_d=20\) units.
\item At \(P=10\): \(Q_d=100-20=80,\; Q_s=20+30=50\).
Shortage \(=Q_d-Q_s=30\) units.
\end{enumerate}

\bigskip

%==========================
% Q2
%==========================
\noindent \textbf{Q2. Shifts in Demand from Income Changes}

Demand: \(Q_d=500-5P+0.5Y\). Supply: \(Q_s=100+3P\).
\begin{enumerate}[label=(\alph*)]
\item If \(Y=200\), find the equilibrium price and quantity.
\item If \(Y=300\), find the new equilibrium price and quantity.
\item Interpret the income effect.
\end{enumerate}

\textbf{Solution.}
\begin{enumerate}[label=(\alph*)]
\item With \(Y=200\): \(Q_d=600-5P\).
\[
600-5P=100+3P \;\eq\; 500=8P \;\eq\; P^\ast=62.5,\quad Q^\ast=100+3(62.5)=287.5.
\]
\item With \(Y=300\): \(Q_d=650-5P\).
\[
650-5P=100+3P \;\eq\; 550=8P \;\eq\; P^\ast=68.75,\quad Q^\ast=100+3(68.75)=306.25.
\]
\item Higher income shifts demand right. Both equilibrium price and quantity rise, indicating laptops behave as a normal good.
\end{enumerate}

\bigskip

%==========================
% Q3
%==========================
\noindent \textbf{Q3. Elasticity at a Point}

For \(Q_d=200-4P\),
\begin{enumerate}[label=(\alph*)]
\item Compute the price elasticity of demand at \(P=30\).
\item Classify demand at this price.
\item Predict the effect on total revenue from a small price increase.
\end{enumerate}

\textbf{Solution.}
\begin{enumerate}[label=(\alph*)]
\item At \(P=30\), \(Q=200-4(30)=80\). For linear demand, \(\varepsilon=\dv{Q}{P}\cdot \frac{P}{Q}=-4\cdot \frac{30}{80}=-1.5\).
\item \(|\varepsilon|=1.5>1\), so demand is elastic.
\item With elastic demand, a price increase lowers total revenue.
\end{enumerate}

\bigskip

%==========================
% Q4
%==========================
\noindent \textbf{Q4. Supply Shift from Input Cost Increase (Graphical and Numerical)}

Initial demand and supply:
\[
Q_d=80-2P,\qquad Q_s=10+2P.
\]
\begin{enumerate}[label=(\alph*))]
\item Find and graph the initial equilibrium.
\item An input cost increase shifts supply to \(Q_s=-10+2P\). Find and graph the new equilibrium.
\item Explain the comparative statics.
\end{enumerate}

\textbf{Solution.}
\begin{enumerate}[label=(\alph*)]
\item Baseline equilibrium:
\[
80-2P=10+2P \;\eq\; 70=4P \;\eq\; P^\ast=17.5,\quad Q^\ast=80-2(17.5)=45.
\]
\item After the shift:
\[
80-2P=-10+2P \;\eq\; 90=4P \;\eq\; P' = 22.5,\quad Q' = 80-2(22.5)=35.
\]
\item Higher input costs shift supply left. Equilibrium price rises and equilibrium quantity falls.
\end{enumerate}

% Optional TikZ sketch (commented): remove % to compile a quick diagram
% \usepackage{tikz,pgfplots} \pgfplotsset{compat=1.18}
% \begin{figure}[h]
% \centering
% \begin{tikzpicture}
% \begin{axis}[width=0.7\textwidth, height=6cm, axis lines=left,
% xlabel={$Q$}, ylabel={$P$}, ymin=0, xmin=0, xmax=90, ymax=40, samples=2, domain=0:90, legend pos=north west]
% \addplot+[mark=none] ({80-2*x}, x); \addlegendentry{$D: Q=80-2P$}
% \addplot+[mark=none] ({10+2*x}, x); \addlegendentry{$S_0: Q=10+2P$}
% \addplot+[mark=none, dashed] ({-10+2*x}, x); \addlegendentry{$S_1: Q=-10+2P$}
% \end{axis}
% \end{tikzpicture}
% \caption{Supply shifts left from $S_0$ to $S_1$.}
% \end{figure}

\bigskip

%==========================
% Q5
%==========================
\noindent \textbf{Q5. Rent Control: Shortage Quantified (Graphical and Numerical)}

Rental market:
\[
Q_d=120-2P,\qquad Q_s=20+2P.
\]
\begin{enumerate}[label=(\alph*)]
\item Find and graph the competitive equilibrium.
\item With a rent ceiling at \(P=20\), show the shortage and compute its size.
\end{enumerate}

\textbf{Solution.}
\begin{enumerate}[label=(\alph*)]
\item Competitive equilibrium:
\[
120-2P=20+2P \;\eq\; 100=4P \;\eq\; P^\ast=25,\quad Q^\ast=120-2(25)=70.
\]
\item At \(P=20\): \(Q_d=120-40=80,\; Q_s=20+40=60\).
Shortage \(=80-60=20\) units. On the graph, plot the horizontal gap between \(Q_d\) and \(Q_s\) at \(P=20\).
\end{enumerate}

% Optional TikZ sketch (commented)
% \begin{figure}[h]
% \centering
% \begin{tikzpicture}
% \begin{axis}[width=0.7\textwidth, height=6cm, axis lines=left,
% xlabel={$Q$}, ylabel={$P$}, ymin=0, xmin=0, xmax=130, ymax=40, samples=2, domain=0:130, legend pos=north east]
% \addplot+[mark=none] ({120-2*x}, x); \addlegendentry{$D$}
% \addplot+[mark=none] ({20+2*x}, x); \addlegendentry{$S$}
% \addplot+[mark=none, dashed] coordinates {(0,20) (130,20)}; \addlegendentry{Ceiling $P=20$}
% \end{axis}
% \end{tikzpicture}
% \caption{Rent ceiling creates a shortage of 20 at $P=20$.}
% \end{figure}

\bigskip

%==========================
% Q6
%==========================
\noindent \textbf{Q6. Comparative Statics of a Demand Shift (Graphical and Numerical)}

Wheat market:
\[
Q_d=600-10P,\qquad Q_s=100+5P.
\]
\begin{enumerate}[label=(\alph*)]
\item Compute the initial equilibrium.
\item Demand shifts right to \(Q_d=700-10P\). Compute the new equilibrium.
\item Explain the movements in price and quantity.
\end{enumerate}

\textbf{Solution.}
\begin{enumerate}[label=(\alph*)]
\item Initial:
\[
600-10P=100+5P \;\eq\; 500=15P \;\eq\; P^\ast=\frac{500}{15}\approx 33.33, \quad Q^\ast=600-10(33.33)\approx 266.67.
\]
\item After shift:
\[
700-10P=100+5P \;\eq\; 600=15P \;\eq\; P' = 40,\quad Q' = 700-10(40)=300.
\]
\item A rightward demand shift raises both equilibrium price and quantity.
\end{enumerate}

\bigskip
\hrule
\bigskip

\textit{Notes for graphing across questions:} Always mark intercepts. For a linear inverse demand \(P=a-bQ\), price intercept \(=a\) and quantity intercept \(=a/b\). Analogous logic applies for supply.

\end{document}
